\documentclass[12pt,a4paper]{article}

\usepackage[T2A]{fontenc}
\usepackage[utf8]{inputenc}
\usepackage[russian]{babel}
\usepackage{geometry}
\usepackage{array}
\usepackage{longtable}
\usepackage{tabularx}
\usepackage{booktabs}
\usepackage{ragged2e}
\usepackage{setspace}

\geometry{left=2cm,right=2cm,top=2cm,bottom=2cm}
\setlength{\parindent}{1.25cm}
\setlength{\parskip}{0.25em}
\renewcommand{\arraystretch}{1.2}
\pagestyle{empty}

\newcolumntype{L}[1]{>{\RaggedRight\arraybackslash}p{#1}}
\newcolumntype{C}[1]{>{\Centering\arraybackslash}p{#1}}

\newcommand{\UniversityName}{Санкт-Петербургский государственный университет}
\newcommand{\DirectionName}{02.04.02 «Фундаментальная информатика и информационные технологии»}
\newcommand{\ProgramName}{Технологии искусственного интеллекта и Big Data}
\newcommand{\StudentName}{Корельский Дмитрий Сергеевич}
\newcommand{\StudentGroup}{\underline{\hspace{2.8cm}}}
\newcommand{\PracticeStart}{01.04.2026}
\newcommand{\PracticeEnd}{15.04.2026}
\newcommand{\PracticePlace}{СПбГУ, факультет прикладной математики -- процессов управления, ОП «Технологии искусственного интеллекта и Big Data»}
\newcommand{\SupervisorName}{Митрофанов Евгений Павлович}
\newcommand{\SupervisorRole}{научный руководитель, СПбГУ}
\newcommand{\PracticeTheme}{Разработка объяснимого прототипа рекомендательной системы выбора следующей культуры в севообороте}

\begin{document}

\begin{center}
{\large ПРИЛОЖЕНИЕ 4}\\[0.5em]
Шаблон титульного листа отчета о практике
\end{center}

\vspace{1.5em}

\begin{center}
\textbf{\MakeUppercase{\UniversityName}}\\[1.5em]
\begin{flushleft}
\textbf{Направление:} \DirectionName\\
\textbf{ОП:} \ProgramName
\end{flushleft}

\vspace{3.5cm}

{\Large\bfseries ОТЧЕТ\\[0.4em]О НАУЧНО-ИССЛЕДОВАТЕЛЬСКОЙ ПРАКТИКЕ}
\end{center}

\vspace{2.5cm}

\noindent\textbf{Тема задания:} \PracticeTheme

\vspace{3cm}

\begin{flushright}
\begin{minipage}{0.58\textwidth}
\textbf{Выполнил:} \StudentName\\
\textit{номер группы:} \StudentGroup

\vspace{1.5em}

\textbf{Руководитель практики:}\\
\SupervisorName\\
\SupervisorRole
\end{minipage}
\end{flushright}

\vfill

\begin{center}
Санкт-Петербург\\
2026
\end{center}

\clearpage
\begin{center}
{\large ПРИЛОЖЕНИЕ 1}\\[0.5em]
Бланк индивидуального задания на практику студента
\end{center}

\vspace{1em}

\begin{center}
\textbf{\MakeUppercase{\UniversityName}}
\end{center}

\noindent\textbf{Направление:} \DirectionName\\
\textbf{ОП:} \ProgramName

\vspace{1.5em}

\begin{center}
{\Large\bfseries ИНДИВИДУАЛЬНОЕ ЗАДАНИЕ}\\[0.3em]
на научно-исследовательскую практику
\end{center}

\vspace{1em}

\noindent\textbf{Студент:} \StudentName \hfill \textbf{номер группы:} \StudentGroup

\noindent\textbf{Тема задания:} \PracticeTheme

\noindent\textbf{Сроки прохождения практики:} \PracticeStart{} -- \PracticeEnd

\noindent\textbf{Место прохождения практики:} \PracticePlace

\vspace{0.75em}

\noindent\textbf{1. Виды работ и требования к их выполнению:}
\begin{enumerate}
\item Провести анализ предметной области и сформулировать задачу как explainable decision-support prototype для рекомендации следующей культуры по истории поля и региональному контексту.
\item Подготовить и проверить воспроизводимый notebook-first pipeline: этапы очистки данных, построения window-датасета и фиксированного split по \texttt{CSBID} с \texttt{random\_state=42}.
\item Исследовать baseline-подходы и обучить финальную CatBoost-модель \texttt{baseline}; сохранить артефакты модели и метаданные в согласованных путях \texttt{artifacts/}.
\item Перейти от top-1 классификации к Top-k recommendation workflow, оценить \texttt{Accuracy@1}, \texttt{Hit@3}, \texttt{Hit@5}, confidence buckets и калибровку вероятностей.
\item Реализовать и оценить explainability-слой на основе knowledge graph и правил поддержки/предупреждения без подмены финальной предиктивной модели.
\item Подготовить spatial demo и итоговые материалы практики: отчет, заполненный дневник, таблицы и графики с опорой на артефакты дипломной работы.
\end{enumerate}

\noindent\textbf{2. Виды отчетных материалов и требования к их оформлению:}
\begin{enumerate}
\item Отчет по практике с описанием цели, постановки задачи, этапов работы, используемых данных, модели \texttt{baseline}, recommendation/confidence/explainability результатов и выводов.
\item Дневник практики с фиксацией ежедневных этапов работ, возникающих вопросов и полученных результатов.
\item Подборка подтверждающих материалов: ноутбуки этапов 01--10, артефакты в \texttt{artifacts/}, итоговые таблицы и рисунки из диплома.
\item При описании результатов использовать только артефактно подтвержденные значения; финальную систему трактовать как прототип поддержки принятия решений, а не как агрономический оптимизатор.
\end{enumerate}

\vspace{1.5em}

\begin{center}
\textbf{3. ПЛАН-ГРАФИК}
\end{center}

\small
\begin{tabularx}{\textwidth}{C{0.8cm}|L{4.2cm}|C{2.2cm}|L{4.5cm}|L{3.1cm}}
\textbf{№ этапа} & \textbf{Наименование этапа} & \textbf{Срок завершения этапа} & \textbf{Виды работ} & \textbf{Форма отчетности} \\
\hline
1 & Постановка задачи и обзор материалов & 03.04.2026 & Анализ темы диплома, структуры репозитория и постановки recommendation-задачи & Рабочие заметки, уточненная тема \\
\hline
2 & Подготовка данных и формирование признаков & 06.04.2026 & Проверка этапов 01--02: очистка данных, window-датасет, feature table & Описание датасета, файлы \texttt{01\_meta.json}, \texttt{02\_meta.json} \\
\hline
3 & Базовые модели и frozen split & 08.04.2026 & Анализ baselines, split по \texttt{CSBID}, контроль контрактов классов & Таблицы baseline-метрик, \texttt{baseline\_meta.json} \\
\hline
4 & Финальная baseline-модель & 10.04.2026 & Анализ CatBoost \texttt{baseline}, feature contract и quality metrics & \texttt{model.cbm}, \texttt{meta.json}, quality plots \\
\hline
5 & Recommendation и explainability слои & 13.04.2026 & Top-k метрики, confidence/calibration, сравнение с Markov, knowledge graph & CSV, PNG и JSON этапов 07--09 \\
\hline
6 & Spatial demo и оформление результатов & 15.04.2026 & Подготовка spatial demo, итогового отчета и приложений практики & Заполненные приложения, итоговый PDF, ссылки на артефакты \\
\end{tabularx}
\normalsize

\vspace{1.5em}

\noindent Задание утверждено на заседании кафедры \underline{\hspace{7cm}}

\noindent (протокол от «\underline{\hspace{0.8cm}}» \underline{\hspace{3cm}} 20\underline{\hspace{0.5cm}} г. №\underline{\hspace{1cm}}).

\vspace{1em}

\noindent Руководитель практики \hfill \underline{\hspace{3cm}} \hfill \SupervisorName

\vspace{1.5em}

\noindent Задание принял к исполнению «\underline{\hspace{0.8cm}}» \underline{\hspace{3cm}} 20\underline{\hspace{0.5cm}} г.

\vspace{1.5em}

\noindent\hfill \underline{\hspace{4cm}} \hfill \StudentName

\clearpage
\begin{center}
{\large ПРИЛОЖЕНИЕ 2}\\[0.5em]
Шаблон дневника практики
\end{center}

\vspace{1em}

\begin{center}
\textbf{\MakeUppercase{\UniversityName}}
\end{center}

\noindent\textbf{Направление:} \DirectionName\\
\textbf{ОП:} \ProgramName

\vspace{1em}

\begin{center}
{\Large\bfseries ДНЕВНИК НАУЧНО-ИССЛЕДОВАТЕЛЬСКОЙ ПРАКТИКИ}
\end{center}

\noindent за период с \PracticeStart{} по \PracticeEnd

\noindent\textbf{Студент:} \StudentName \hfill \textbf{номер группы:} \StudentGroup

\noindent\textbf{Руководитель практики:} \SupervisorName, \SupervisorRole

\vspace{1em}

\scriptsize
\begin{longtable}{L{1.55cm}|L{2.55cm}|L{4.6cm}|L{2.25cm}|L{3.0cm}|L{1.55cm}}
\textbf{Дата} & \textbf{Наименование структурного подразделения} & \textbf{Краткое содержание работы} & \textbf{Возникшие вопросы} & \textbf{Достигнутые результаты} & \textbf{Отметка о выполнении} \\
\hline
01.04.2026 & СПбГУ, ПМ-ПУ & Уточнена тема НИР и прикладная постановка: recommendation system для next-crop выбора, explainable prototype, top-k workflow. & Как корректно ограничить claims рамкой decision support. & Зафиксирована финальная framing без завышения выводов. & Выполнено \\
\hline
02.04.2026 & СПбГУ, ПМ-ПУ & Проверены результаты этапа 01: очистка и агрегация данных USDA NASS CSB/CDL, подготовленный датасет на 8\,110\,870 строк. & Какие классы исключать как нерелевантные. & Подтвержден артефакт \texttt{01\_df\_prepared\_filtered.pkl}. & Выполнено \\
\hline
03.04.2026 & СПбГУ, ПМ-ПУ & Проанализирован этап 02: формирование window-датасета, исторических признаков и target-таблицы. & Как сохранить совместимость признаков для downstream notebooks. & Подтвержден window-набор на 38\,511\,210 строк и 11 столбцов. & Выполнено \\
\hline
06.04.2026 & СПбГУ, ПМ-ПУ & Проверен frozen split по \texttt{CSBID} и baseline-эксперименты этапа 03. & Как избежать утечки между train/val/test. & Подтвержден group-aware split и файл \texttt{baseline\_meta.json}. & Выполнено \\
\hline
07.04.2026 & СПбГУ, ПМ-ПУ & Изучена финальная baseline CatBoost-модель и ее feature contract: \texttt{history\_1..3}, \texttt{STATEFIPS}, \texttt{CNTYFIPS}, \texttt{CSBACRES}, \texttt{INSIDE\_X}, \texttt{INSIDE\_Y}. & Достаточен ли baseline как финальная модель диплома. & Подтвержден статус \texttt{baseline} как final model в \texttt{model\_registry.json}. & Выполнено \\
\hline
08.04.2026 & СПбГУ, ПМ-ПУ & Проанализированы recommendation-метрики CatBoost в Top-k режиме. & Как интерпретировать качество за пределами top-1. & Для test подтверждены \texttt{Acc@1=0.699}, \texttt{Hit@3=0.951}, \texttt{Hit@5=0.991}. & Выполнено \\
\hline
09.04.2026 & СПбГУ, ПМ-ПУ & Изучены confidence buckets и калибровка вероятностей на этапе 07. & Насколько вероятности пригодны как confidence-сигнал. & Подтверждены ECE test \texttt{0.017} и high-confidence \texttt{Acc@1=0.841}. & Выполнено \\
\hline
10.04.2026 & СПбГУ, ПМ-ПУ & Проведено сравнение с Markov-1 recommendation baseline и анализ практической полезности shortlist-подхода. & Насколько простой последовательностный baseline конкурентоспособен. & Зафиксировано преимущество CatBoost над Markov-1 по всем Top-k метрикам. & Выполнено \\
\hline
13.04.2026 & СПбГУ, ПМ-ПУ & Проанализирован explainability-слой: knowledge graph, правила поддержки/предупреждения, case studies. & Как не смешивать интерпретацию и предсказание. & Подтверждены 2\,827 узлов, 18\,154 ребра и 5 правил в KG-слое. & Выполнено \\
\hline
14.04.2026 & СПбГУ, ПМ-ПУ & Рассмотрен spatial demo: карта с полями, рекомендациями и explainability-контекстом. & Как показать результаты на карте без переобучения модели. & Подтверждены HTML-артефакты spatial demo и bbox-демо. & Выполнено \\
\hline
15.04.2026 & СПбГУ, ПМ-ПУ & Подготовлены материалы отчета по практике и приложений на основе фактических артефактов дипломной работы. & Какие поля требуют ручного уточнения перед сдачей. & Сформирован комплект приложений 1, 2 и 4 в \LaTeX{}-виде. & Выполнено \\
\end{longtable}
\normalsize

\vspace{0.5em}

\noindent * Подпись руководителя практики от организации


\end{document}
